\documentclass[twocolumn]{aastex631}
\begin{document}
\title{The reionization ionizing-photon-budget shortfall for star-forming galaxies is not robust to systematics at z\~{}5 (highC systematic corner)}
\author{NebulaMind Lab (autonomous pipeline)}
\affiliation{NebulaMind Lab, \url{https://lab.nebulamind.net}}
\begin{abstract}
Reionization ionizing-photon-budget reconciliation at z\~{}5: star-forming galaxies require f\_esc=0.029 (+0.027/-0.014) to close the budget (Madau-Dickinson SFRD, log xi\_ion=25.5+/-0.15, clumping C=2-5, JWST-SFRD tail), versus indirect-proxy-inferred f\_esc=0.062 (+0.108/-0.039) from LzLCS O32/beta calibrations. Median delta(required-inferred)=-0.031 dex-frac (16-84\%: -0.139 to +0.012); 25\% of the systematic MC shows a shortfall. Conclusion: the budget CLOSES within the systematic, and the sign holds under both O32 and beta calibrations. The result is bounded by the xi\_ion x clumping x proxy-calibration systematic, not by statistics. Generated autonomously from public data (jwst) via the NebulaMind Lab runner: a bounded, reproducible, descriptive study.
\end{abstract}
\section{Introduction}
Recent studies have highlighted a potential crisis in the photon budget for reionization, with some suggesting that star-forming galaxies may not produce enough ionizing photons to account for the observed reionization [Muñoz2024]. This has led to increased scrutiny of the assumptions and calibrations used in estimating the ionizing photon budget. In particular, there is a need to reconcile the cosmic star formation rate density (SFRD) with the escape fraction of ionizing photons from galaxies, as well as the clumping factor of the intergalactic medium [Park2022].
\section{Data and method}
To address this issue, we have performed a literature-anchored budget calculation using established values for the SFRD, ionization efficiency (xi\_ion), and escape fraction (f\_esc) proxy calibrations. Specifically, we adopt the Madau \& Dickinson (2014) analytic fitting function for the cosmic SFRD, and use published values for xi\_ion and f\_esc from LzLCS [Chisholm+22, Flury+22; Simmonds+24]. We do not utilize any survey catalog data or observational results from JWST, SDSS, or TNG in this analysis.
\section{Result}
Our calculations indicate that star-forming galaxies require an escape fraction of f\_esc = 0.029 (+0.027/-0.014) to close the reionization photon budget at z\~{}5, assuming a clumping factor C between 2 and 5 [Madau2017]. This value is lower than the indirect-proxy-inferred escape fraction of f\_esc = 0.062 (+0.108/-0.039) derived from LzLCS O32/beta calibrations. The median difference between the required and inferred escape fractions is -0.031 dex-frac, with a range of -0.139 to +0.012 (16-84\% confidence interval). Notably, 25\% of our systematic Monte Carlo simulations show a shortfall in the photon budget.
\begin{figure}\centering\includegraphics[width=\columnwidth]{result.png}\caption{The reionization ionizing-photon-budget shortfall for star-forming galaxies is not robust to systematics at z\~{}5 (highC systematic corner)}\end{figure}
\section{Caveats}
It is important to note that this result is subject to several caveats and limitations. Firstly, our analysis relies on a single selection of literature values for key parameters, which may not fully capture the underlying uncertainties and variations in these quantities. Secondly, we have not accounted for potential systematic errors or biases in the calibrations used to infer f\_esc from observational data. Finally, our calculation assumes a fixed clumping factor C between 2 and 5, which may not accurately reflect the true complexity of the intergalactic medium during reionization. These limitations highlight the need for further research and improved constraints on the key parameters governing the ionizing photon budget.
\section*{References}
\footnotesize
[Muoz2024] Muñoz et al. (2024). Reionization after JWST: a photon budget crisis?. bibcode 2024MNRAS.535L..37M \par [Park2022] Park et al. (2022). Calibrating excursion set reionization models to approximately conserve ionizing photons. bibcode 2022MNRAS.517..192P \par [Duncan2015] Duncan et al. (2015). Powering reionization: assessing the galaxy ionizing photon budget at z < 10. bibcode 2015MNRAS.451.2030D \par [Davies2021] Davies et al. (2021). The Predicament of Absorption-dominated Reionization: Increased Demands on Ionizing Sources. bibcode 2021ApJ...918L..35D \par [Madau2017] Madau (2017). Cosmic Reionization after Planck and before JWST: An Analytic Approach. bibcode 2017ApJ...851...50M

\end{document}
